
%(BEGIN_QUESTION)
% Copyright 2011, Tony R. Kuphaldt, released under the Creative Commons Attribution License (v 1.0)
% This means you may do almost anything with this work of mine, so long as you give me proper credit

Anta at en radiosender gir ut 327 watt RF-effekt ved frekvens 1070 kHz, inn i en halv-bølge ``whip''-antenne.  Radioenergien stråles ut fra sendeantennen og mottas av en annen antenne, slik at signalet inn i mottakeren blir 9.5 mikrowatt.

Beregn senderens uteffekt ($P_{tx}$) i både dBm og dBW, og mottakerens inngangssignal ($P_{rx}$) i dBm og dBW.  Beregn deretter totalt effekttap mellom sender og mottaker ($P_{loss}$), i dB.

\vskip 10pt

$P_{tx}$ = \underbar{\hskip 50pt} dBm = \underbar{\hskip 50pt} dBW

\vskip 10pt

$P_{rx}$ = \underbar{\hskip 50pt} dBm = \underbar{\hskip 50pt} dBW

\vskip 10pt

$P_{loss}$ = \underbar{\hskip 50pt} dB

\vskip 20pt \vbox{\hrule \hbox{\strut \vrule{} {\bf Forslag til sokratisk diskusjon} \vrule} \hrule}

\begin{itemize}
\item{} Forklar hvorfor ``decibel'' er en så vanlig enhet for effekt, og for effektforsterkning/tap, i RF-kommunikasjon.
\item{} Estimer disse decibelverdiene ved først å runde effektforholdet til nærmeste produkt av 10 og/eller 2, og bruk så ekvivalensen 10 dB for 10-gangers forhold og 3 dB for 2-gangers forhold.
\end{itemize}

\underbar{file i00285}
%(END_QUESTION)





%(BEGIN_ANSWER)

$P_{tx}$ = \underbar{55.15} dBm = \underbar{25.15} dBW

\vskip 10pt

$P_{rx}$ = \underbar{$-20.22$} dBm = \underbar{$-50.22$} dBW

\vskip 10pt

$P_{loss}$ = \underbar{$-75.37$} dB

%(END_ANSWER)





%(BEGIN_NOTES)


%INDEX% Electronics review: decibel power calculations

%(END_NOTES)
